% ═══════════════════════════════════════════════════════════════════════════
%  APPENDICE DE LIAISON — identique dans les sept articles.
%  % ═══════════════════════════════════════════════════════════════════════════
%  APPENDICE DE LIAISON — identique dans les sept articles.
%  % ═══════════════════════════════════════════════════════════════════════════
%  APPENDICE DE LIAISON — identique dans les sept articles.
%  % ═══════════════════════════════════════════════════════════════════════════
%  APPENDICE DE LIAISON — identique dans les sept articles.
%  \input{appendix-where} juste avant \begin{thebibliography} ou la biblio.
%  Il répond à : qui vit où, et qu'est-ce que la série a changé en cours de route.
% ═══════════════════════════════════════════════════════════════════════════
\section*{Appendix: which sector lives where, and what changed}
\addcontentsline{toc}{section}{Appendix: which sector lives where}

The series was written over several months and \textbf{two of its assignments moved}.
A reader who meets the articles out of order is owed an explicit map rather than an
inference, so it is given here, identically in all seven papers.

\subsection*{A.1 The three locations}
The geometry has three places where physics can sit: the \emph{brane} at $y=0$,
the \emph{bulk} of the $8.2\,\mu$m interval, and the \emph{far wall} at $y=\pi R$.
\begin{center}\small
\begin{tabular}{llll}\toprule
sector & lives on & carried by & established in\\\midrule
ordinary matter & \textbf{brane} & thin core, Standard Model & I\\
gravity, radion, graviton & \textbf{bulk} & five-dimensional fields & I\\
the mediator $\Phi$ and its tower & \textbf{bulk} & Airy states of the $|y|$ well & II\\
the galactic response & \textbf{bulk} $\to$ brane & $\Phi$ exchange between baryons & III\\
dark matter & \textbf{see A.2} & --- & III, revisited in VI\\
dark energy & \textbf{see A.3} & --- & III, revisited in VI\\
hidden gauge sector & \textbf{far wall} & eight D8 and one O8 & V\\\bottomrule
\end{tabular}
\end{center}

\subsection*{A.2 Dark matter: the reservoir, and what the wall reading adds}
Articles~I--IV place dark matter in the \emph{bulk}: the excited states of the Airy
tower, populated cold at birth, collisionless, carrying at least $98.6\%$ of the dark
sector's density. That assignment does the work in three places. It makes structure
formation identical to cold dark matter at linear order. It answers the Bullet Cluster
by mass budget alone --- cluster masses are cluster masses, and the lensing tracks the
reservoir, which is why the lensing--dynamics difficulty of purely modified dynamics is
not inherited. And it fixes the radiation budget of Article~I.

Article~VI proposes a different reading: dark matter as the ordinary matter of the six
other regions of the same wall, seen through the bulk, giving
$\Omega_{\rm DM}/\Omega_b=N_{\rm regions}-1=6$ against an observed $5.31$.
\textbf{These are not the same statement, and we do not pretend otherwise.}

What they share is everything the phenomenology uses: in both readings the dark
component is cold, collisionless, and carries the mass, so lensing, the Bullet Cluster
and linear structure formation are unaffected. What differs is the \emph{constituent}:
quanta of $\Phi$ in the first, baryons of other regions in the second.
\textbf{The two are compatible if the tower reservoir is not required to be the whole
of the dark matter} --- and Article~III's own argument only requires it to dominate the
budget, not to be exhaustive. \Open

\subsection*{A.3 Dark energy: a bulk scalar, then a brane curvature}
Articles~I--IV treat $\Phi$ as the dark-energy candidate --- a bulk scalar engine. That
attempt \emph{failed by forty orders of magnitude}: a four-dimensional scalar cannot
reach $2.4$~meV from a $10^{-13}$~eV engine scale. The failure is recorded, not hidden.

Article~VI relocates dark energy \emph{to the brane}, as the extrinsic curvature of the
D8 wall in the cosmological bulk, $\Lambda_4=\varepsilon^2(1/R)^4$ with
$\varepsilon\simeq10^{-2}$. The mechanism supplies $(1/R)^4$ outright and leaves a gap
of $10^4$ rather than $10^{40}$. \textbf{$\Phi$ is correspondingly demoted from
dark-energy candidate to dark-sector mediator}, which is the role Articles~II and III
already gave it in every equation they actually use. \Candidate, declared tuning.

\subsection*{A.4 The radial acceleration relation, and what makes galaxies turn}
The relation is the one place where brane and bulk are \emph{both} indispensable, and
where nothing has moved since Article~III.

Baryons sit on the brane. They source $\Phi$, which propagates in the bulk with mass
$m_\Phi=24$~meV and hence Compton range $8.2\,\mu$m --- microscopic, so no fifth force
acts between laboratory bodies. What acts at galactic scale is not the free field but
the \emph{medium}: the coherent component of the dark sector responds to the baryonic
distribution, and that response, read back on the brane, is the extra acceleration.
The response saturates, which is why it produces a relation --- an acceleration scale
$g_\dagger$ --- rather than a rescaling of Newton's constant.

Three consequences, all in Article~III and unchanged:
\begin{itemize}\itemsep2pt
\item the response couples to the trace $T^\mu_{\ \mu}$, hence to mass and not to
      composition; the universality is bounded on 171 SPARC galaxies;
\item the response is carried by at most $1.4\%$ of the dark sector, the coherent part,
      so it never competes with the reservoir in the mass budget --- \textbf{galaxies
      turn because of the medium, they weigh because of the reservoir};
\item the response \emph{dies} in weak field, which is the ultra-diffuse galaxy test,
      and it dies in a decoherence bubble around the Sun, which is why Cassini is
      satisfied identically rather than by tuning.
\end{itemize}
\textbf{Neither relocation of A.2 or A.3 touches this.} The galactic response was never
assigned to the dark-energy sector, and it is not the reservoir either; it is the third
role, and it is the one the framework was built to explain.

\subsection*{A.5 What a reader should carry between articles}
Articles~I--IV stand on their own evidence and are not modified by V--VII: no result of
theirs is derived from the string embedding, and the embedding cannot promote any of
their labels. What V--VII add is an ultraviolet origin for one constant, a candidate
geometry, and two consequences of that geometry. What they \emph{change} is recorded
above and nowhere else: two relocations, both declared, neither retroactive.
 juste avant \begin{thebibliography} ou la biblio.
%  Il répond à : qui vit où, et qu'est-ce que la série a changé en cours de route.
% ═══════════════════════════════════════════════════════════════════════════
\section*{Appendix: which sector lives where, and what changed}
\addcontentsline{toc}{section}{Appendix: which sector lives where}

The series was written over several months and \textbf{two of its assignments moved}.
A reader who meets the articles out of order is owed an explicit map rather than an
inference, so it is given here, identically in all seven papers.

\subsection*{A.1 The three locations}
The geometry has three places where physics can sit: the \emph{brane} at $y=0$,
the \emph{bulk} of the $8.2\,\mu$m interval, and the \emph{far wall} at $y=\pi R$.
\begin{center}\small
\begin{tabular}{llll}\toprule
sector & lives on & carried by & established in\\\midrule
ordinary matter & \textbf{brane} & thin core, Standard Model & I\\
gravity, radion, graviton & \textbf{bulk} & five-dimensional fields & I\\
the mediator $\Phi$ and its tower & \textbf{bulk} & Airy states of the $|y|$ well & II\\
the galactic response & \textbf{bulk} $\to$ brane & $\Phi$ exchange between baryons & III\\
dark matter & \textbf{see A.2} & --- & III, revisited in VI\\
dark energy & \textbf{see A.3} & --- & III, revisited in VI\\
hidden gauge sector & \textbf{far wall} & eight D8 and one O8 & V\\\bottomrule
\end{tabular}
\end{center}

\subsection*{A.2 Dark matter: the reservoir, and what the wall reading adds}
Articles~I--IV place dark matter in the \emph{bulk}: the excited states of the Airy
tower, populated cold at birth, collisionless, carrying at least $98.6\%$ of the dark
sector's density. That assignment does the work in three places. It makes structure
formation identical to cold dark matter at linear order. It answers the Bullet Cluster
by mass budget alone --- cluster masses are cluster masses, and the lensing tracks the
reservoir, which is why the lensing--dynamics difficulty of purely modified dynamics is
not inherited. And it fixes the radiation budget of Article~I.

Article~VI proposes a different reading: dark matter as the ordinary matter of the six
other regions of the same wall, seen through the bulk, giving
$\Omega_{\rm DM}/\Omega_b=N_{\rm regions}-1=6$ against an observed $5.31$.
\textbf{These are not the same statement, and we do not pretend otherwise.}

What they share is everything the phenomenology uses: in both readings the dark
component is cold, collisionless, and carries the mass, so lensing, the Bullet Cluster
and linear structure formation are unaffected. What differs is the \emph{constituent}:
quanta of $\Phi$ in the first, baryons of other regions in the second.
\textbf{The two are compatible if the tower reservoir is not required to be the whole
of the dark matter} --- and Article~III's own argument only requires it to dominate the
budget, not to be exhaustive. \Open

\subsection*{A.3 Dark energy: a bulk scalar, then a brane curvature}
Articles~I--IV treat $\Phi$ as the dark-energy candidate --- a bulk scalar engine. That
attempt \emph{failed by forty orders of magnitude}: a four-dimensional scalar cannot
reach $2.4$~meV from a $10^{-13}$~eV engine scale. The failure is recorded, not hidden.

Article~VI relocates dark energy \emph{to the brane}, as the extrinsic curvature of the
D8 wall in the cosmological bulk, $\Lambda_4=\varepsilon^2(1/R)^4$ with
$\varepsilon\simeq10^{-2}$. The mechanism supplies $(1/R)^4$ outright and leaves a gap
of $10^4$ rather than $10^{40}$. \textbf{$\Phi$ is correspondingly demoted from
dark-energy candidate to dark-sector mediator}, which is the role Articles~II and III
already gave it in every equation they actually use. \Candidate, declared tuning.

\subsection*{A.4 The radial acceleration relation, and what makes galaxies turn}
The relation is the one place where brane and bulk are \emph{both} indispensable, and
where nothing has moved since Article~III.

Baryons sit on the brane. They source $\Phi$, which propagates in the bulk with mass
$m_\Phi=24$~meV and hence Compton range $8.2\,\mu$m --- microscopic, so no fifth force
acts between laboratory bodies. What acts at galactic scale is not the free field but
the \emph{medium}: the coherent component of the dark sector responds to the baryonic
distribution, and that response, read back on the brane, is the extra acceleration.
The response saturates, which is why it produces a relation --- an acceleration scale
$g_\dagger$ --- rather than a rescaling of Newton's constant.

Three consequences, all in Article~III and unchanged:
\begin{itemize}\itemsep2pt
\item the response couples to the trace $T^\mu_{\ \mu}$, hence to mass and not to
      composition; the universality is bounded on 171 SPARC galaxies;
\item the response is carried by at most $1.4\%$ of the dark sector, the coherent part,
      so it never competes with the reservoir in the mass budget --- \textbf{galaxies
      turn because of the medium, they weigh because of the reservoir};
\item the response \emph{dies} in weak field, which is the ultra-diffuse galaxy test,
      and it dies in a decoherence bubble around the Sun, which is why Cassini is
      satisfied identically rather than by tuning.
\end{itemize}
\textbf{Neither relocation of A.2 or A.3 touches this.} The galactic response was never
assigned to the dark-energy sector, and it is not the reservoir either; it is the third
role, and it is the one the framework was built to explain.

\subsection*{A.5 What a reader should carry between articles}
Articles~I--IV stand on their own evidence and are not modified by V--VII: no result of
theirs is derived from the string embedding, and the embedding cannot promote any of
their labels. What V--VII add is an ultraviolet origin for one constant, a candidate
geometry, and two consequences of that geometry. What they \emph{change} is recorded
above and nowhere else: two relocations, both declared, neither retroactive.
 juste avant \begin{thebibliography} ou la biblio.
%  Il répond à : qui vit où, et qu'est-ce que la série a changé en cours de route.
% ═══════════════════════════════════════════════════════════════════════════
\section*{Appendix: which sector lives where, and what changed}
\addcontentsline{toc}{section}{Appendix: which sector lives where}

The series was written over several months and \textbf{two of its assignments moved}.
A reader who meets the articles out of order is owed an explicit map rather than an
inference, so it is given here, identically in all seven papers.

\subsection*{A.1 The three locations}
The geometry has three places where physics can sit: the \emph{brane} at $y=0$,
the \emph{bulk} of the $8.2\,\mu$m interval, and the \emph{far wall} at $y=\pi R$.
\begin{center}\small
\begin{tabular}{llll}\toprule
sector & lives on & carried by & established in\\\midrule
ordinary matter & \textbf{brane} & thin core, Standard Model & I\\
gravity, radion, graviton & \textbf{bulk} & five-dimensional fields & I\\
the mediator $\Phi$ and its tower & \textbf{bulk} & Airy states of the $|y|$ well & II\\
the galactic response & \textbf{bulk} $\to$ brane & $\Phi$ exchange between baryons & III\\
dark matter & \textbf{see A.2} & --- & III, revisited in VI\\
dark energy & \textbf{see A.3} & --- & III, revisited in VI\\
hidden gauge sector & \textbf{far wall} & eight D8 and one O8 & V\\\bottomrule
\end{tabular}
\end{center}

\subsection*{A.2 Dark matter: the reservoir, and what the wall reading adds}
Articles~I--IV place dark matter in the \emph{bulk}: the excited states of the Airy
tower, populated cold at birth, collisionless, carrying at least $98.6\%$ of the dark
sector's density. That assignment does the work in three places. It makes structure
formation identical to cold dark matter at linear order. It answers the Bullet Cluster
by mass budget alone --- cluster masses are cluster masses, and the lensing tracks the
reservoir, which is why the lensing--dynamics difficulty of purely modified dynamics is
not inherited. And it fixes the radiation budget of Article~I.

Article~VI proposes a different reading: dark matter as the ordinary matter of the six
other regions of the same wall, seen through the bulk, giving
$\Omega_{\rm DM}/\Omega_b=N_{\rm regions}-1=6$ against an observed $5.31$.
\textbf{These are not the same statement, and we do not pretend otherwise.}

What they share is everything the phenomenology uses: in both readings the dark
component is cold, collisionless, and carries the mass, so lensing, the Bullet Cluster
and linear structure formation are unaffected. What differs is the \emph{constituent}:
quanta of $\Phi$ in the first, baryons of other regions in the second.
\textbf{The two are compatible if the tower reservoir is not required to be the whole
of the dark matter} --- and Article~III's own argument only requires it to dominate the
budget, not to be exhaustive. \Open

\subsection*{A.3 Dark energy: a bulk scalar, then a brane curvature}
Articles~I--IV treat $\Phi$ as the dark-energy candidate --- a bulk scalar engine. That
attempt \emph{failed by forty orders of magnitude}: a four-dimensional scalar cannot
reach $2.4$~meV from a $10^{-13}$~eV engine scale. The failure is recorded, not hidden.

Article~VI relocates dark energy \emph{to the brane}, as the extrinsic curvature of the
D8 wall in the cosmological bulk, $\Lambda_4=\varepsilon^2(1/R)^4$ with
$\varepsilon\simeq10^{-2}$. The mechanism supplies $(1/R)^4$ outright and leaves a gap
of $10^4$ rather than $10^{40}$. \textbf{$\Phi$ is correspondingly demoted from
dark-energy candidate to dark-sector mediator}, which is the role Articles~II and III
already gave it in every equation they actually use. \Candidate, declared tuning.

\subsection*{A.4 The radial acceleration relation, and what makes galaxies turn}
The relation is the one place where brane and bulk are \emph{both} indispensable, and
where nothing has moved since Article~III.

Baryons sit on the brane. They source $\Phi$, which propagates in the bulk with mass
$m_\Phi=24$~meV and hence Compton range $8.2\,\mu$m --- microscopic, so no fifth force
acts between laboratory bodies. What acts at galactic scale is not the free field but
the \emph{medium}: the coherent component of the dark sector responds to the baryonic
distribution, and that response, read back on the brane, is the extra acceleration.
The response saturates, which is why it produces a relation --- an acceleration scale
$g_\dagger$ --- rather than a rescaling of Newton's constant.

Three consequences, all in Article~III and unchanged:
\begin{itemize}\itemsep2pt
\item the response couples to the trace $T^\mu_{\ \mu}$, hence to mass and not to
      composition; the universality is bounded on 171 SPARC galaxies;
\item the response is carried by at most $1.4\%$ of the dark sector, the coherent part,
      so it never competes with the reservoir in the mass budget --- \textbf{galaxies
      turn because of the medium, they weigh because of the reservoir};
\item the response \emph{dies} in weak field, which is the ultra-diffuse galaxy test,
      and it dies in a decoherence bubble around the Sun, which is why Cassini is
      satisfied identically rather than by tuning.
\end{itemize}
\textbf{Neither relocation of A.2 or A.3 touches this.} The galactic response was never
assigned to the dark-energy sector, and it is not the reservoir either; it is the third
role, and it is the one the framework was built to explain.

\subsection*{A.5 What a reader should carry between articles}
Articles~I--IV stand on their own evidence and are not modified by V--VII: no result of
theirs is derived from the string embedding, and the embedding cannot promote any of
their labels. What V--VII add is an ultraviolet origin for one constant, a candidate
geometry, and two consequences of that geometry. What they \emph{change} is recorded
above and nowhere else: two relocations, both declared, neither retroactive.
 juste avant \begin{thebibliography} ou la biblio.
%  Il répond à : qui vit où, et qu'est-ce que la série a changé en cours de route.
% ═══════════════════════════════════════════════════════════════════════════
\section*{Appendix: which sector lives where, and what changed}
\addcontentsline{toc}{section}{Appendix: which sector lives where}

The series was written over several months and \textbf{two of its assignments moved}.
A reader who meets the articles out of order is owed an explicit map rather than an
inference, so it is given here, identically in all seven papers.

\subsection*{A.1 The three locations}
The geometry has three places where physics can sit: the \emph{brane} at $y=0$,
the \emph{bulk} of the $8.2\,\mu$m interval, and the \emph{far wall} at $y=\pi R$.
\begin{center}\small
\begin{tabular}{llll}\toprule
sector & lives on & carried by & established in\\\midrule
ordinary matter & \textbf{brane} & thin core, Standard Model & I\\
gravity, radion, graviton & \textbf{bulk} & five-dimensional fields & I\\
the mediator $\Phi$ and its tower & \textbf{bulk} & Airy states of the $|y|$ well & II\\
the galactic response & \textbf{bulk} $\to$ brane & $\Phi$ exchange between baryons & III\\
dark matter & \textbf{see A.2} & --- & III, revisited in VI\\
dark energy & \textbf{see A.3} & --- & III, revisited in VI\\
hidden gauge sector & \textbf{far wall} & eight D8 and one O8 & V\\\bottomrule
\end{tabular}
\end{center}

\subsection*{A.2 Dark matter: the reservoir, and what the wall reading adds}
Articles~I--IV place dark matter in the \emph{bulk}: the excited states of the Airy
tower, populated cold at birth, collisionless, carrying at least $98.6\%$ of the dark
sector's density. That assignment does the work in three places. It makes structure
formation identical to cold dark matter at linear order. It answers the Bullet Cluster
by mass budget alone --- cluster masses are cluster masses, and the lensing tracks the
reservoir, which is why the lensing--dynamics difficulty of purely modified dynamics is
not inherited. And it fixes the radiation budget of Article~I.

Article~VI proposes a different reading: dark matter as the ordinary matter of the six
other regions of the same wall, seen through the bulk, giving
$\Omega_{\rm DM}/\Omega_b=N_{\rm regions}-1=6$ against an observed $5.31$.
\textbf{These are not the same statement, and we do not pretend otherwise.}

What they share is everything the phenomenology uses: in both readings the dark
component is cold, collisionless, and carries the mass, so lensing, the Bullet Cluster
and linear structure formation are unaffected. What differs is the \emph{constituent}:
quanta of $\Phi$ in the first, baryons of other regions in the second.
\textbf{The two are compatible if the tower reservoir is not required to be the whole
of the dark matter} --- and Article~III's own argument only requires it to dominate the
budget, not to be exhaustive. \Open

\subsection*{A.3 Dark energy: a bulk scalar, then a brane curvature}
Articles~I--IV treat $\Phi$ as the dark-energy candidate --- a bulk scalar engine. That
attempt \emph{failed by forty orders of magnitude}: a four-dimensional scalar cannot
reach $2.4$~meV from a $10^{-13}$~eV engine scale. The failure is recorded, not hidden.

Article~VI relocates dark energy \emph{to the brane}, as the extrinsic curvature of the
D8 wall in the cosmological bulk, $\Lambda_4=\varepsilon^2(1/R)^4$ with
$\varepsilon\simeq10^{-2}$. The mechanism supplies $(1/R)^4$ outright and leaves a gap
of $10^4$ rather than $10^{40}$. \textbf{$\Phi$ is correspondingly demoted from
dark-energy candidate to dark-sector mediator}, which is the role Articles~II and III
already gave it in every equation they actually use. \Candidate, declared tuning.

\subsection*{A.4 The radial acceleration relation, and what makes galaxies turn}
The relation is the one place where brane and bulk are \emph{both} indispensable, and
where nothing has moved since Article~III.

Baryons sit on the brane. They source $\Phi$, which propagates in the bulk with mass
$m_\Phi=24$~meV and hence Compton range $8.2\,\mu$m --- microscopic, so no fifth force
acts between laboratory bodies. What acts at galactic scale is not the free field but
the \emph{medium}: the coherent component of the dark sector responds to the baryonic
distribution, and that response, read back on the brane, is the extra acceleration.
The response saturates, which is why it produces a relation --- an acceleration scale
$g_\dagger$ --- rather than a rescaling of Newton's constant.

Three consequences, all in Article~III and unchanged:
\begin{itemize}\itemsep2pt
\item the response couples to the trace $T^\mu_{\ \mu}$, hence to mass and not to
      composition; the universality is bounded on 171 SPARC galaxies;
\item the response is carried by at most $1.4\%$ of the dark sector, the coherent part,
      so it never competes with the reservoir in the mass budget --- \textbf{galaxies
      turn because of the medium, they weigh because of the reservoir};
\item the response \emph{dies} in weak field, which is the ultra-diffuse galaxy test,
      and it dies in a decoherence bubble around the Sun, which is why Cassini is
      satisfied identically rather than by tuning.
\end{itemize}
\textbf{Neither relocation of A.2 or A.3 touches this.} The galactic response was never
assigned to the dark-energy sector, and it is not the reservoir either; it is the third
role, and it is the one the framework was built to explain.

\subsection*{A.5 What a reader should carry between articles}
Articles~I--IV stand on their own evidence and are not modified by V--VII: no result of
theirs is derived from the string embedding, and the embedding cannot promote any of
their labels. What V--VII add is an ultraviolet origin for one constant, a candidate
geometry, and two consequences of that geometry. What they \emph{change} is recorded
above and nowhere else: two relocations, both declared, neither retroactive.
